\section{Free entropy dimensions and von Neumann generators}
\label{sec:entropy-dimensions}

For a bounded self-adjoint tuple $Z$ in a tracial von Neumann algebra,
let $\delta(Z)$ and $\delta_0(Z)$ denote Voiculescu's microstates free
entropy dimensions. The second is the modified version using entropy
in the presence of the perturbing semicircular family. Write $\delta^*(Z)$
and $\delta^\star(Z)$ for the nonmicrostates dimensions defined through
nonmicrostates free entropy $\chi^*$ and free Fisher information
$\Phi^*$, respectively. We use the conventions in
\cite[Section~2.4]{voiculescu2002entropy} and
\cite[Section~4.1.1]{connes_shlyakhtenko2005}.

By their definitions, all four quantities depend only on the joint
tracial $*$-distribution of the tuple. This is weaker than invariance
under a change of von Neumann generators: two tuples can generate the
same von Neumann algebra without having the same joint distribution.
The isomorphisms proved above show that the four dimensions do change
under such a replacement.

\begin{corollary}[Dependence on von Neumann generators]
\label{cor:entropy-dimensions}
Let $M=L(\F_2)$ with its canonical trace. For every integer $n\geq2$
there are a self-adjoint $n$-tuple $X^{(n)}$ and a self-adjoint
$2n$-tuple $Y^{(n)}$ such that
\[
 W^*(X^{(n)})=W^*(Y^{(n)})=M
\]
and
\[
 \delta(X^{(n)})=\delta_0(X^{(n)})=n,
 \qquad
 \delta^*(Y^{(n)})=\delta^\star(Y^{(n)})=n.
\]
Consequently, none of $\delta,\delta_0,\delta^*,\delta^\star$ is
invariant under changing a finite self-adjoint $W^*$-generating tuple
of a tracial von Neumann algebra. Each takes unbounded values on such
generating tuples of the single factor $L(\F_2)$.
\end{corollary}

\begin{proof}
Fix an integer $n\geq2$. Theorem~\ref{thm:main} and successive
applications of Theorem~\ref{thm:rank-step} provide a unital normal
trace-preserving $*$-isomorphism
\[
 \theta_n:L(\F_n)\longrightarrow M;
\]
for $n=2$ we may take the identity. If $Z$ is any finite self-adjoint tuple in
$L(\F_n)$ and $p$ is a noncommutative $*$-polynomial, then
\[
 \tau_M\bigl(p(\theta_n(Z))\bigr)=\tau_{L(\F_n)}\bigl(p(Z)\bigr).
\]
Thus coordinatewise transport by $\theta_n$ preserves the joint
tracial $*$-distribution and hence each of the four dimensions.
Normality and surjectivity also give
$W^*(\theta_n(Z))=\theta_n(W^*(Z))$.

Let $S^{(n)}=(s_1,\ldots,s_n)$ be a freely independent variance-one semicircular
tuple in the semicircular realization of $L(\F_n)$
\cite[Remarks~1.10--1.11]{voiculescu1990}. It generates $L(\F_n)$ as a von Neumann algebra,
and the standard semicircular computations give
\[
 \delta(S^{(n)})=\delta_0(S^{(n)})=n;
\]
see \cite[Section~2.4(c) and (d)]{voiculescu2002entropy} for $\delta$
and \cite[Section~5, opening paragraph]{brannan_elzinga_harris_yamashita2023}
for $\delta_0$. Set $X^{(n)}=\theta_n(S^{(n)})$. The preceding
transport observations prove the required generation and the two
microstates values.

For the nonmicrostates dimensions, let $u_1,\ldots,u_n$ be the
canonical group unitaries in $L(\F_n)$ and put
\[
 a_j=\frac{u_j+u_j^*}{2},\qquad
 b_j=\frac{u_j-u_j^*}{2\ii},\qquad
 G^{(n)}=(a_1,b_1,\ldots,a_n,b_n).
\]
The entries are self-adjoint, and $u_j=a_j+\ii b_j$ and
$u_j^*=a_j-\ii b_j$. Hence they generate the algebraic group algebra
$\C\F_n$ as a unital $*$-algebra, and generate $L(\F_n)$ as a von
Neumann algebra. Apply Mineyev and Shlyakhtenko's group-algebra theorem
\cite[Theorem~4.1]{mineyev_shlyakhtenko2005} to this tuple, before
transporting it. It gives
\[
 \delta^*(G^{(n)})=\delta^\star(G^{(n)})
 =\beta_1^{(2)}(\F_n)-\beta_0^{(2)}(\F_n)+1=n.
\]
Here $\beta_0^{(2)}(\F_n)=0$ and $\beta_1^{(2)}(\F_n)=n-1$ are the
usual $L^2$-Betti numbers of the free group; see
\cite[Section~2.1.1]{connes_shlyakhtenko2005}. Define
$Y^{(n)}=\theta_n(G^{(n)})$. Again its joint tracial $*$-distribution
is unchanged and it generates $M$, proving the remaining assertions.
Taking $n=2$ and $n=3$ gives different values on generators of the
same factor, and letting $n$ grow gives unboundedness.
\end{proof}

The nonmicrostates computation uses algebraic generators of
$\C\F_n$ before transport. The transported tuples all generate
$L(\F_2)$ as a von Neumann algebra, but their algebraic $*$-algebras
are not held fixed as $n$ varies. Thus the corollary concerns
$W^*$-generation and is compatible with the algebraic-generator
invariance in \cite[Theorem~4.1]{mineyev_shlyakhtenko2005}.
